\begin{python}
from scripts.calc_1 import *
\end{python}

\section{Нормирование параметров и переменных цепи}

В соответствии с рекомендациями, для упрощения дальнейших расчетов производится нормирование параметров элементов цепи. В качестве базисной частоты принимается значение $\omega_{\text{б}} = \pv[.0e]{rec['omega_base']} \text{ с}^{-1}$, в качестве базисного сопротивления --- сопротивление резистора $R_1$, то есть $R_{\text{б}} = \pv[.0]{rec['R_base']} \text{ Ом}$.
 
\begin{figure}[H]
    \centering
    \begin{tikzpicture}[transform shape]
        % Paths, nodes and wires:
        \draw (1, 1) to[american current source, l={$i_1$}] (1, 4);
        \draw (3, 4) to[european resistor, l_={$R_1$}] (3, 1);
        \draw (8, 4) to[european resistor, l={$R_2$}] (8, 1);
        \draw (1, 1) -- (3, 1) -- (5, 1) -- (8, 1);
        \draw (1, 4) -- (3, 4) -- (5, 4);
        \node[inputarrow, rotate=-90] at (3, 3.25){};
        \node[inputarrow, rotate=-90] at (8, 3.25){};
        \node[circ](N1) at (1, 4){} node[anchor=south east] at (N1.north west){$1$};
        \node[circ](N2) at (1, 1){} node[anchor=north east] at (N2.south west){$4$};
        \node[circ](N3) at (8, 4){} node[anchor=south west] at (N3.north east){$2$};
        \node[circ](N4) at (5, 2.375){} node[anchor=west] at (N4.east){$3$};
        \draw (5, 2.25) to[capacitor, l={$C_1$}] (5, 1);
        \draw (5, 4) to[american inductor, l={$L_1$}] (5, 2.25);
        \draw (5, 4) to[american inductor, l={$L_2$}] (8, 4);
        \node[shape=rectangle, minimum width=-0.035cm, minimum height=-0.035cm](N5) at (3.5, 3){} node[anchor=center] at (N5.text){$+$} node[anchor=north west, align=left, text width=-0.423cm, inner sep=6pt] at (3.5, 3){};
        \node[shape=rectangle, minimum width=-0.035cm, minimum height=-0.035cm](N6) at (3.5, 2){} node[anchor=center] at (N6.text){$-$} node[anchor=north west, align=left, text width=-0.423cm, inner sep=6pt] at (3.5, 2){};
        \node[shape=rectangle, minimum width=-0.035cm, minimum height=-0.035cm](N7) at (7.5, 3){} node[anchor=center] at (N7.text){$+$} node[anchor=north west, align=left, text width=-0.423cm, inner sep=6pt] at (7.5, 3){};
        \node[shape=rectangle, minimum width=-0.035cm, minimum height=-0.035cm](N8) at (7.5, 2){} node[anchor=center] at (N8.text){$-$} node[anchor=north west, align=left, text width=-0.423cm, inner sep=6pt] at (7.5, 2){};
        \node[inputarrow] at (5.75, 4){};
        \node[inputarrow, rotate=-90] at (5, 3.75){};
        \node[inputarrow, rotate=-90] at (5, 2){};
        \node[shape=rectangle, minimum width=0.215cm, minimum height=0.215cm](N9) at (6, 3.75){} node[anchor=center] at (N9.text){$+$} node[anchor=north west, align=left, text width=-0.173cm, inner sep=6pt] at (5.875, 3.875){};
        \node[shape=rectangle, minimum width=0.244cm, minimum height=0.244cm](N10) at (7, 3.75){} node[anchor=center] at (N10.text){$-$} node[anchor=north west, align=left, text width=-0.144cm, inner sep=6pt] at (6.86, 3.89){};
        \node[shape=rectangle, minimum width=0.311cm, minimum height=0.369cm](N11) at (4.75, 3.702){} node[anchor=center] at (N11.text){$+$} node[anchor=north west, align=left, text width=-0.077cm, inner sep=6pt] at (4.577, 3.904){};
        \node[shape=rectangle, minimum width=0.325cm, minimum height=0.325cm](N12) at (4.75, 2.57){} node[anchor=center] at (N12.text){$-$} node[anchor=north west, align=left, text width=-0.063cm, inner sep=6pt] at (4.57, 2.75){};
        \node[shape=rectangle, minimum width=0.311cm, minimum height=0.369cm](N13) at (4.32, 1.813){} node[anchor=center] at (N13.text){$+$} node[anchor=north west, align=left, text width=-0.077cm, inner sep=6pt] at (4.147, 2.015){};
        \node[shape=rectangle, minimum width=0.325cm, minimum height=0.325cm](N14) at (4.32, 1.43){} node[anchor=center] at (N14.text){$-$} node[anchor=north west, align=left, text width=-0.063cm, inner sep=6pt] at (4.14, 1.61){};
    \end{tikzpicture}
    \caption{Исходная схема электрической цепи}
    \label{fig:scheme_orig}
\end{figure}

Формулы перехода к нормированным значениям и результаты вычислений:

\begin{equation}
    R_{1*} = \frac{R_1}{R_{\text{б}}} = \frac{\pv[.0]{s['R1']}}{\pv[.0]{rec['R_base']}} = \pv[.4]{n['R1']}
\end{equation}

\begin{equation}
    R_{2*} = \frac{R_2}{R_{\text{б}}} = \frac{\pv[.0]{s['R2']}}{\pv[.0]{rec['R_base']}} = \pv[.4]{n['R2']}
\end{equation}

\begin{equation}
    L_{1*} = \frac{L_1 \omega_{\text{б}}}{R_{\text{б}}} = \frac{\pv[.4]{s['L1']} \cdot \pv[.0]{rec['omega_base']}}{\pv[.0]{rec['R_base']}} = \pv[.4]{n['L1']}
\end{equation}

\begin{equation}
    L_{2*} = \frac{L_2 \omega_{\text{б}}}{R_{\text{б}}} = \frac{\pv[.3]{s['L2']} \cdot \pv[.0]{rec['omega_base']}}{\pv[.0]{rec['R_base']}} \approx \pv[.4]{n['L2']}
\end{equation}

\begin{equation}
    C_{1*} = C_1 \cdot R_{\text{б}} \cdot \omega_{\text{б}} = \pv[.0e]{s['C1']} \cdot \pv[.0]{rec['R_base']} \cdot \pv[.0e]{rec['omega_base']} \approx \pv[.4]{n['C1']}
\end{equation}

Нормирование временных параметров:

\begin{equation}
    t_{\text{и}1*} = t_{\text{и}1} \omega_{\text{б}} = \pv[.2e]{s['ti1']} \cdot \pv[.0e]{rec['omega_base']} = \pv[.1]{n['ti1']}
\end{equation}

\begin{equation}
    t_{\text{и}2*} = t_{\text{и}2} \omega_{\text{б}} = \pv[.2e]{s['ti2']} \cdot \pv[.0e]{rec['omega_base']} = \pv[.2]{n['ti2']}
\end{equation}

Далее индексы нормирования «$*$» опускаются для упрощения записи. Все последующие расчеты ведутся в нормированных величинах.